\chapter{Introduction}
\label{ch:intro}
A central task in financial engineering is to build models that reproduce observed market
options. A good model must do two things. It must be complex enough to match the observed
cross-section of option prices, in particular the implied volatility. It must also be tractable
enough so that prices can be computed quickly and reliably, since calibration requires doing so.

Affine jump-diffusions (AJDs) are a class of models that a balance between the two.
In an AJD, the drift, the instantaneous covariance, the jump intensity and the short rate are all
affine functions of the state vector. This structure allows for an surprising amount of tractability
while still allowing for stochastic volatility, jumps in prices and jumps in volatility.

The key result is the analysis of \parencite{Duffie2000}, that shows that the extended transform
\begin{equation}
  \phi^{\chi}(v,u,X_t,t,T)
  = \mathbb{E}\!\left[ \ee^{-\int_t^T R(X_s)\,\mathrm{d}s}\,(v_0+v_1\cdot X_T)\,\ee^{u\cdot X_T} \,\middle|\, \mathcal{F}_t \right]
  \label{eq:extended-transform}
\end{equation}
is exponential-affine in the state variables, with coefficients that solve a system of ordinary
differential equations (ODEs), for some models the ODEs even have closed-form solutions.

The extended transform is a powerful tool, by considering special cases of this transform one can
obtain different quantities, e.g:
\begin{itemize}
  \item With $R = v_1 = 0, v_1 = 1$ $\phi^{\chi}$ becomes the we get characteristic function of $X$.
  \item Putting $v_0 = 1, v_1 = u =0$ yields the risk-free zero coupon bond price with maturity $T$.
\end{itemize}

This project studies this modelling approach from theory to data. We first present the AJD framework
and the affine transform formula of \parencite{Duffie2000}. We then derive the \parencite{Lewis2001}
pricing formula, which turns the characteristic function into a European call price through a single
integral. As an example of AJD models we study the particularly tractable SVJJ model, in which the
log-price and the variance jump simultaneously. We derive the models jump transform and the
present the closed-form solution of its Riccati equations. Finally, we calibrate SVJJ to S\&P~500
index options observed on 13 November 2019 and compare it with two simpler benchmarks: the SVJ
model, which has no variance jumps, and the \parencite{Heston1993} model, which has no jumps at
all.

The rest of the project is organised as follows. \Cref{ch:affine} introduces AJD models, the
extended transform and Fourier pricing. \Cref{ch:svjj} presents the SVJJ model.
\Cref{ch:calibration} describes the data and filtering procedure, the calibration setup and
the results. \Cref{ch:conclusion} concludes. Additional results and details are available in the
appendices.
